
\section*{CHAPTER OVERVIEW}
\section*{1: Linear Equations}
\section*{Learning Objectives}
In this chapter, you will learn to:

\begin{enumerate}
  \item Graph a linear equation.
  \item Find the slope of a line.
  \item Determine an equation of a line.
  \item Solve linear systems.
  \item Do application problems using linear equations.
\end{enumerate}

\section*{1: Graphing a Linear Equation}
1.1.1: Graphing a Linear Equation (Exercises)\\
1.2: Slope of a Line\\
1.2.1: Slope of a Line (Exercises)\\
1.3: Determining the Equation of a Line\\
1.3.1: Determining the Equation of a Line (Exercises)\\
1.4: Applications\\
1.4.1: Applications (Exercises)\\
1.5: More Applications\\
1.5.1: More Applications (Exercises)

\section*{6: Chapter Review}
Thumbnail: The red and blue lines on this graph have the same slope (gradient); the red and green lines have the same $y$-intercept (cross the $y$-axis at the same place). (CC BY-SA 1.0; ElectroKid via Wikipedia)

This page titled 1: Linear Equations is shared under a CC BY 4.0 license and was authored, remixed, and/or curated by Rupinder Sekhon and Roberta Bloom via source content that was edited to the style and standards of the LibreTexts platform.

\section*{1: Graphing a Linear Equation}
\section*{Learning Objectives}
In this section, you will leam to:

\begin{enumerate}
  \item Graph a line when you know its equation
  \item Graph a line when you are given its equation in parametric form
  \item Graph and find equations of vertical and horizontal lines
\end{enumerate}

\section*{Graphing a Line from it Equation}
Equations whose graphs are straight lines are called linear equations. The following are some examples of linear equations:\\
$2 x-3 y=6, \quad 3 x=4 y-7, \quad y=2 x-5, \quad 2 y=3, \quad$ and $\quad x-2=0$\\
A line is completely determined by two points. Therefore, to graph a linear equation we need to find the coordinates of two points. This can be accomplished by choosing an arbitrary value for x or y and then solving for the other variable.

\section*{Example 1.1.1}
Graph the line: $y=3 x+2$

\section*{Solution}
We need to find the coordinates of at least two points. We arbitrarily choose $\mathrm{x}=-1, \mathrm{x}=0$, and $\mathrm{x}=1$.

\begin{itemize}
  \item If $x=-1$, then $y=3(-1)+2$ or -1 . Therefore, $(-1,-1)$ is a point on this line.
  \item If $x=0$, then $y=3(0)+2$ or $y=2$. Hence the point $(0,2)$.
  \item If $x=1$, then $y=5$, and we get the point $(1,5)$.
\end{itemize}

Below, the results are summarized, and the line is graphed.

\begin{center}
\begin{tabular}{|l|l|l|l|}
\hline
$x$ & -1 & 0 & 1 \\
\hline
$y$ & -1 & 2 & 5 \\
\hline
\end{tabular}
\end{center}

\begin{center}
\includegraphics[width=\textwidth]{2024_12_11_69ed6e5e21c55666356eg-02}
\end{center}

\section*{Example 1.1.2}
Graph the line: $2 x+y=4$

\section*{Solution}
Again, we need to find coordinates of at least two points.\\
We arbitrarily choose $x=-1, x=0$, and $y=2$.

\begin{itemize}
  \item If $x=-1$, then $2(-1)+y=4$ which results in $y=6$. Therefore, $(-1,6)$ is a point on this line.
  \item If $x=0$, then $2(0)+y=4$, which results in $y=4$. Hence the point $(0,4)$.
  \item If $y=2$, then $2 x+2=4$, which yields $x=1$, and gives the point $(1,2)$.
\end{itemize}

The table below shows the points, and the line is graphed.

\begin{center}
\begin{tabular}{|l|l|l|l|l|l|l|l|l|}
\hline
$x$ & -1 & 0 & 1 \\
\hline
$y$ & 6 & 4 & 2 \\
\hline
\end{tabular}
\end{center}

\begin{center}
\includegraphics[width=\textwidth]{2024_12_11_69ed6e5e21c55666356eg-03}
\end{center}

\section*{Intercepts}
The points at which a line crosses the coordinate axes are called the intercepts.\\
When graphing a line by plotting two points, using the intercepts is often preferred because they are easy to find.

\begin{itemize}
  \item To find the value of the $x$-intercept, we let $y=0$
  \item To find the value of the $y$-intercept, we let $x=0$.
\end{itemize}

\section*{Example 1.1.3}
Find the intercepts of the line: $2 x-3 y=6$, and graph.

\section*{Solution}
To find the x -intercept, let $\mathrm{y}=0$ in the equation, and solve for x .

$$
\begin{array}{r}
2 x-3(0)=6 \\
2 x-0=6 \\
2 x=6 \\
x=3
\end{array}
$$

Therefore, the $x$-intercept is the point $(3,0)$.\\
To find the y -intercept, let $\mathrm{x}=0$ in the equation, and solve for y .

$$
\begin{aligned}
2(0)-3 y & =6 \\
0-3 y & =6 \\
-3 y & =6 \\
y & =-2
\end{aligned}
$$

Therefore, the $y$-intercept is the point $(0,-2)$.\\
To graph the line, plot the points for the $x$-intercept $(3,0)$ and the $y$-intercept $(0,-2)$, and use them to draw the line.\\
\includegraphics[max width=\textwidth, center]{2024_12_11_69ed6e5e21c55666356eg-03(1)}

\section*{Graphing a Line from its Equation in Parametric Form}
In higher math, equations of lines are sometimes written in parametric form. For example, $x=3+2 t, y=1+t$. The letter $t$ is called the parameter, or the dummy variable.\\
Parametric lines can be graphed by finding values for x and y by substituting numerical values for t . Plot the points using their ( $\mathrm{x}, \mathrm{y}$ ) coordinates and use the points to draw the line.

\section*{Example 1.1.4}
Graph the line given by the parametric equations: $x=3+2 t, y=1+t$

\section*{Solution}
Let $\mathrm{t}=0,1$ and 2 ; for each value of t find the corresponding values for x and y .\\
The results are given in the table below.

\begin{center}
\begin{tabular}{|c|c|c|c|}
\hline
$\mathbf{t}$ & $\mathbf{x}$ & $\mathbf{y}$ & Point co Line \\
\hline
0 & 3 & 1 & $(3,1)$ \\
\hline
1 & 5 & 2 & $(5,2)$ \\
\hline
2 & 7 & 3 & $(7,3)$ \\
\hline
\end{tabular}
\end{center}

\begin{center}
\includegraphics[width=\textwidth]{2024_12_11_69ed6e5e21c55666356eg-04}
\end{center}

\section*{Horizontal and Vertical Lines}
When an equation of a line has only one variable, the resulting graph is a horizontal or a vertical line.

\begin{itemize}
  \item The graph of the line $x=a$, where $a$ is a constant, is a vertical line that passes through the point $(a, 0)$. Every point on this line has the $x$-coordinate equal to $a$, regardless of the $y$-coordinate.
\end{itemize}

\begin{itemize}
  \item The graph of the line $y=b$, where $b$ is a constant, is a horizontal line that passes through the point $(0, b)$. Every point on this line has the $y$-coordinate equal to $b$, regardless of the $x$-coordinate.
\end{itemize}

\section*{Example 1.1.5}
Graph the lines: $x=-2$, and $y=3$.

\section*{Solution}
The graph of the line $x=-2$ is a vertical line that has the $x$-coordinate -2 no matter what the $y$-coordinate is. The graph is a vertical line passing through point $(-2,0)$.\\
The graph of the line $y=3$, is a horizontal line that has the $y$-coordinate 3 regardless of what the $x$-coordinate is. Therefore, the graph is a horizontal line that passes through point $(0,3)$.\\
\includegraphics[max width=\textwidth, center]{2024_12_11_69ed6e5e21c55666356eg-05}

Note: Most students feel that the coordinates of points must always be integers. This is not true, and in real life situations, not always possible. Do not be intimidated if your points include numbers that are fractions or decimals.

This page titled 1.1: Graphing a Linear Equation is shared under a CC BY 4.0 license and was authored, remixed, and/or curated by Rupinder Sekhon and Roberta Bloom via source contemt that was edited to the style and standards of the LibreTexts platform.

\subsection*{1: Graphing a Linear Equation (Exercises)}
\section*{SECTION 1.1 PROBLEM SET: GRAPHING A LINEAR EQUATION}
Work the following problems.

\begin{enumerate}
  \item Is the point $(2,3)$ on the line $5 x-2 y=4$ ?
\end{enumerate}

\begin{center}
\begin{tabular}{ll}
3) For the line $3 x-y=12$, complete the following ordered pairs. & 2) Is the point $(1,-2)$ on the line $6 x-y=4$ ? \\
\begin{tabular}{ll}
(2, $)(, 6)$ & (3, $)(, 4)$ \\
(0, $)(, 0)$ & (0, $)(, 0)$ \\
5) Graph $y=2 x+3$ & 6) Graph $y=-3 x+5$ \\
7) Graph $y=4 x-3$ & 8) Graph $x-2 y=8$ \\
9) Graph $2 x+y=4$ & 10) Graph $2 x-3 y=6$ \\
\end{tabular} &  \\
\end{tabular}
\end{center}

\section*{SECTION 1.1 PROBLEM SET: GRAPHING A LINEAR EQUATION}
\begin{enumerate}
  \setcounter{enumi}{10}
  \item Graph $2 x+4=0$
  \item Graph the following three equations on the same set of courdinate axes.
\end{enumerate}

$$
y=x+1 \quad y=2 x+1 \quad y=-x+1
$$

\begin{enumerate}
  \setcounter{enumi}{14}
  \item Graph the line using the parametric equations $\mathrm{x}=1+2 \mathrm{t}, \mathrm{y}=3+\mathrm{t}$
  \item Greph $2 y-6=0$
  \item Graph the following three equations on the same set of coordinate axes.
\end{enumerate}

$$
y=2 x+1 \quad y=2 x \quad y=2 x-1
$$

\begin{enumerate}
  \setcounter{enumi}{15}
  \item Graph the line using the parametric equations $x=2$ - 3t, $y=1+2 t$
\end{enumerate}

This page titled 1.1.1: Graphing a Linear Equation (Exercises) is shared under a CC BY 4.0 license and was authored, remixed, and/or curated by Rupinder Sekhon and Roberta Bloom via source content that was edited to the style and standards of the LibreTexts platform.

\section*{2: Slope of a Line}
\section*{Learning Objectives}
In this section, you will leam to:

\begin{enumerate}
  \item Find the slope of a line.
  \item Graph the line if a point and the slope are given.
\end{enumerate}

In the last section, we learned to graph a line by choosing two points on the line. A graph of a line can also be determined if one point and the "steepness" of the line is known. The number that refers to the steepness or inclination of a line is called the slope of the line.

From previous math courses, many of you remember slope as the "rise over run," or "the vertical change over the horizontal change" and have often seen it expressed as:

$$
\frac{\text { rise }}{\text { run }}, \frac{\text { vertical change }}{\text { horizontal change }}, \frac{\Delta y}{\Delta x} \text { etc. }
$$

We give a precise definition.

\section*{Definition: Slope}
If $\left(x_{1}, y_{1}\right)$ and $\left(x_{2}, y_{2}\right)$ are two different points on a line, the slope of the line is


\begin{equation*}
\text { slope }=m=\frac{y_{2}-y_{1}}{x_{2}-x_{1}} \tag{1.2.1}
\end{equation*}


\section*{Example 1.2.1}
Find the slope of the line passing through points $(-2,3)$ and $(4,-1)$, and graph the line.

\section*{Solution}
Let $\left(x_{1}, y_{1}\right)=(-2,3)$ and $\left(x_{2}, y_{2}\right)=(4,-1)$, then the slope (via Equation 1.2.1) is

$$
m=\frac{-1-3}{4-(-2)}=\frac{-4}{6}=-\frac{2}{3}
$$

\begin{center}
\includegraphics[width=\textwidth]{2024_12_11_69ed6e5e21c55666356eg-07}
\end{center}

To give the reader a better understanding, both the vertical change, -4 , and the horizontal change, 6 , are shown in the above figure.

When two points are given, it does not matter which point is denoted as ( $x_{1}, y_{1}$ ) and which ( $x_{2}, y_{2}$ ). The value for the slope will be the same.\\
In Example 1.2.1, If we instead choose $\left(x_{1}, y_{1}\right)=(4,-1)$ and $\left(x_{2}, y_{2}\right)=(-2,3)$, then we will get the same value for the slope as we obtained earlles.\\
The steps involved are as follows.

$$
m=\frac{3-(-1)}{-2-4}=\frac{4}{-6}=-\frac{2}{3}
$$

The student should further observe that

\begin{itemize}
  \item If a line rises when going from left to right, then it has a positive slope. In this situation, as the value of $x$ increases, the value of $y$ also increases
  \item If a line falls going from left to right, it has a negative slope; as the value of $x$ increases, the value of $y$ decreases.
\end{itemize}

\section*{Example 1.2.2}
Find the slope of the line that passes through the points $(2,3)$ and $(2,-1)$, and graph.

\section*{Solution}
Let $\left(x_{1}, y_{1}\right)=(2,3)$ and $\left(x_{2}, y_{2}\right)=(2,-1)$ then the slope is

$$
m=\frac{-1-3}{2-2}=\frac{4}{0}=\text { undefined }
$$

\begin{center}
\includegraphics[width=\textwidth]{2024_12_11_69ed6e5e21c55666356eg-08(1)}
\end{center}

Note: The slope of a vertical line is undefined.

\section*{Example 1.2.3}
Find the slope of the line that passes through the points ( $-1,-4$ ) and (3, 4 )

\section*{Solution}
Let $\left(x_{1}, y_{1}\right)=(-1,-4)$ and $\left(x_{2}, y_{2}\right)=(3,-4)$, then the slope is

$$
m=\frac{-4-(-4)}{3-(-1)}=\frac{0}{4}=0
$$

\begin{center}
\includegraphics[width=\textwidth]{2024_12_11_69ed6e5e21c55666356eg-08}
\end{center}

Note: The slope of a horizontal line is 0

\section*{Example 1.2.4}
Graph the line that passes through the point $(1,2)$ and has slope $-\frac{3}{4}$.

\section*{Solution}
Slope equals $\frac{\mathrm{rise}}{\mathrm{rum}}$. The fact that the slope is $\frac{-3}{4}$, means that for every rise of -3 units (fall of 3 units) there is a run of 4 . So if from the given point ( 1,2 ) we go down 3 units and go right 4 units, we reach the point ( $5,-1$ ). The graph is obtained by connecting these two points.\\
\includegraphics[max width=\textwidth, center]{2024_12_11_69ed6e5e21c55666356eg-09(1)}

Alternatively, since $\frac{3}{-4}$ represents the same number, the line can be drawn by starting at the point $(1,2)$ and choosing a rise of 3 units followed by a run of -4 units. So from the point $(1,2)$, we go up 3 units, and to the left 4 , thus reaching the point $(-3,5)$ which is also on the same line. See figure below.\\
\includegraphics[max width=\textwidth, center]{2024_12_11_69ed6e5e21c55666356eg-09}

\section*{Example 1.2.5}
Find the slope of the line $2 x+3 y=6$.

\section*{Solution}
In order to find the slope of this line, we will choose any two points on this line.\\
Again, the selection of $x$ and $y$ intercepts seems to be a good choice. The $x$-intercept is (3, 0), and the $y$-intercept is ( 0,2 ). Therefore, the slope is

$$
m=\frac{2-0}{0-3}=-\frac{2}{3}
$$

The graph below shows the line and the $x$-intercepts and $y$-intercepts:\\
\includegraphics[max width=\textwidth, center]{2024_12_11_69ed6e5e21c55666356eg-09(2)}

\section*{Example 1.2.6}
Find the slope of the line $y=3 x+2$.

\section*{Solution}
We again find two points on the line, e.g., $(0,2)$ and $(1,5)$. Therefore, the slope is

$$
m=\frac{5-2}{1-0}=\frac{3}{1}=3
$$

Look at the slopes and the $y$-intercepts of the following lines.

\begin{center}
\begin{tabular}{|c|c|c|}
\hline
The line & slope & $y$-intercept \\
\hline
$y=3 z+2$ & 3 & 2 \\
\hline
$y=-2 x+5$ & -2 & 5 \\
\hline
$y=\frac{3}{2} x-4$ & $\frac{2}{2}$ & -4 \\
\hline
\end{tabular}
\end{center}

It is no coincidence that when an equation of the line is solved for $y$, the coefficient of the $x$ term represents the slope, and the constant term represents the $y$-intercept.\\
In other words, for the line $y=m x+b, m$ is the slope, and $b$ is the $y$-intercept.

\section*{Example 1.2.7}
Determine the slope and $y$-intercept of the line $2 x+3 y=6$.

\section*{Solution}
We solve for $y$ :

$$
\begin{aligned}
& 2 x+3 y=6 \\
& 3 y=-2 x+6 \\
& y=(-2 / 3) x+2
\end{aligned}
$$

The slope $=$ the coefficient of the $x$ term $=-2 / 3$.\\
The $y$-intercept $=$ the constant term $=2$.

This page titled 1.2: Slope of a Line is shared under a CC BY 4.0 license and was authored, remixed, and/or curated by Rupinder Sekhon and Roberta Bloom via source content that was edited to the style and standards of the LibreTexts platform.

\subsection*{1: Slope of a Line (Exercises)}
\section*{SECTION 1.2 PROBLEM SET: SLOPE OF A LINE}
Find the slope of the line passing through the following pair of points.

\begin{center}
\begin{tabular}{|l|l|}
\hline
1) $(2,3)$ and $(5,9)$ & 2) $(4,1)$ and $(2,5)$ \\
\hline
3) $(-1,1)$ and $(1,3)$ & 4) $(4,3)$ and $(-1,3)$ \\
\hline
5) $(6,-5)$ and $(4,-1)$ & 6) $(5,3)$ and $(-1,-4)$ \\
\hline
7$)(3,4)$ and $(3,7)$ & $8)(-2,4)$ and $(-3,-2)$ \\
\hline
9$)(-3,-5)$ and $(-1,-7)$ & $10)(0,4)$ and $(3,0)$ \\
\hline
\end{tabular}
\end{center}

\section*{SECTION 1.2 PROBLEM SET: SLOPE OF A LINE}
Determine the slope of the line from the given equation of the line.\\
11) $y=-2 x+1$\\
12) $y=3 x-2$\\
13) $2 x-y=6$\\
15) $3 x-4 y=12$\\
14) $x+3 y=6$\\
16) What is the slope of the $x$-axis? What is the slope of the $y$-axis?

Graph the line that passes through the given point and has the given slope.\\
17) $(1,2)$ and $m=-3 / 4$\\
18) $(2,-1)$ and $m=2 / 3$\\
19) $(0,2)$ and $m=-2$\\
20) $(2,3)$ and $m=0$

This page titled 1.2.1: Slope of a Line (Exercises) is shared under a CC BY 4.0 license and was authored, remixed, and/or curated by Rupinder Sekhon and Roberta Bloom via source content that was edited to the style and standards of the LibreTexts platform.

\section*{3: Determining the Equation of a Line}
\section*{Learning Objectives}
In this section, you will leam to:

\begin{enumerate}
  \item Find an equation of a line if a point and the slope are given.
  \item Find an equation of a line if two points are given.
\end{enumerate}

So far, we were given an equation of a line and were asked to give information about it. For example, we were asked to find points on the line, find its slope and even find intercepts. Now we are going to reverse the process. That is, we will be given either two points, or a point and the slope of a line, and we will be asked to find its equation.\\
An equation of a line can be written in three forms, the slope-intercept form, the point-slope form, or the standard form. We will discuss each of them in this section.

A line is completely determined by two points, or by a point and slope. The information we are given about a particular line will influence which form of the equation is most convenient to use. Once we know any form of the equation of a line, it is easy to reexpress the equation in the other forms if needed.

\section*{THE SLOPE-INTERCEPT FORM OF A LINE: $\boldsymbol{y}=\boldsymbol{m} \boldsymbol{x}+\boldsymbol{b}$}
In the last section we learned that the equation of a line whose slope $=m$ and $y$-intercept $=b$ is

$$
\mathbf{y}=\mathbf{m x}+\mathbf{b}
$$

This is called the slope-intercept form of the line and is the most commonly used form.

\section*{Example 1.3.1}
Find an equation of a line whose slope is 5 , and $y$-intercept is 3 .

\section*{Solution}
Since the slope is $m=5$, and the $y$-intercept is $b=3$, the equation is $y=5 x+3$.

\section*{Example 1.3.2}
Find the equation of the line that passes through the point $(2,7)$ and has slope 3 .

\section*{Solution}
Since $m=3$, the partial equation is $y=3 x+b$.\\
Now $b$ can be determined by substituting the point $(2,7)$ in the equation $y=3 x+b$.

$$
\begin{aligned}
& 7=3(2)+b \\
& b=1
\end{aligned}
$$

Therefore, the equation is $y=3 x+1$.

\section*{Example 1.3.3}
Find an equation of the line that passes through the points $(-1,2)$, and ( 1,8 ).

\section*{Solution}
$m=\frac{8-2}{1-(-1)}=\frac{6}{2}=3$. So the partial equation is $y=3 x+b$.\\
We can use elther of the two points $(-1,2)$ or $(1,8)$, to find $b$. Substituting $(-1,2)$ gives

$$
\begin{aligned}
& 2=3(-1)+b \\
& 5=b
\end{aligned}
$$

So the equation is $y=3 x+5$.

\section*{Example 1.3.4}
Find an equation of the line that has $x$-intercept 3 , and $y$-intercept 4 .

\section*{Solution}
$x$-intercept $=3$, and $y$-intercept $=4$ correspond to the points (3, 0), and ( 0,4 ), respectively.

$$
m=\frac{4-0}{0-3}=-\frac{4}{3}
$$

We are told the $y$-intercept is 4 ; thus $b=4$\\
Therefore, the equation is $y=-\frac{4}{3} x+4$.

\section*{THE POINT-SLOPE FORM OF A LINE: $y-y_{1}=\boldsymbol{m}\left(\boldsymbol{x}-\boldsymbol{x}_{1}\right)$}
The point-slope form is useful when we know two points on the line and want to find the equation of the line.\\
Let $L$ be a line with slope $m$, and known to contain a specific point ( $x_{1}, y_{1}$ ). If ( $x, y$ ) is any other point on the line $L$, then the definition of a slope leads us to the point-slope form or point-slope formula.

The slope is $\frac{z-y_{1}}{x-x_{1}}=m$\\
Multiplying both sides by ( $x-x_{1}$ ) gives the point-slope form:

$$
\mathrm{y}-\mathrm{y}_{1}=\mathrm{m}\left(\mathrm{x}-\mathrm{x}_{1}\right)
$$

\section*{Example 1.3.5}
Find the point-slope form of the equation of a line that has slope 1.5 and passes through the point $(12,4)$.

\section*{Solution}
Substituting the point $\left(x_{1}, y_{1}\right)=(12,4)$ and $m=1.5$ in the point-slope formula, we get

$$
\begin{gathered}
y-y_{1}=\mathbf{m}\left(x-x_{1}\right) \\
y-4=1.5(x-12)
\end{gathered}
$$

The student may be tempted to simplify this into the slope intercept form $y=m x+b$. But since the problem specifically requests point-slope form we will not simplify it.

\section*{THE STANDARD FORM OF A LINE: $\boldsymbol{A} \boldsymbol{x}+\boldsymbol{B} \boldsymbol{y}=\boldsymbol{C}$}
Another useful form of the equation of a line is the standard form.\\
If we know the equation of a line in point-slope form, $y-y_{1}=m\left(x-x_{1}\right)$, or if we know the equation of the line in slopeintercept form $y=m x+b$, we can simplify the formula to have all terms for the $x$ and $y$ variables on one side of the equation, and the constant on the other side of the equation.\\
The result is referred to as the standard form of the line:

$$
\mathbf{A x}+\mathbf{B y}=\mathbf{C}
$$

\section*{Example 1.3.6}
Using the point-slope formula, find the standard form of an equation of the line that passes through the point ( 2,3 ) and has slope $-3 / 5$.\\
Solution: Substituting the point $(2,3)$ and $m=-3 / 5$ in the point-slope formula, we get

$$
y-3=-3 / 5(x-2)
$$

Multiplying both sides by 5 gives us

$$
\begin{aligned}
& 5(y-3)=-3(x-2) \\
& 5 y-15=-3 x+6 \\
& 3 x+5 y=21 \quad \text { Standard Form }
\end{aligned}
$$

\section*{Example 1.3.7}
Find the standard form of the line that passes through the points $(1,-2)$, and $(4,0)$.

\section*{Solution}
First we find the slope: $m=\frac{0-(-2)}{4-1}=\frac{2}{3}$\\
Then, the point-slope form is: $y-(-2)=\frac{2}{3}(x-1)$\\
Multiplying both sides by 3 gives us

$$
\begin{aligned}
& 3(y+2)=2(x-1) \\
& 3 y+6=2 x-2 \\
& -2 x+3 y=-8 \\
& 2 x-3 y=8 \quad \text { Standard Form }
\end{aligned}
$$

We should always be able to convert from one form of an equation to anothes. For example, if we are given a line in the slopeintercept form, we should be able to express it in the standard form, and vice versa.

\section*{Example 1.3.8}
Write the equation $y=-\frac{2}{3} x+3$ in the standard form.

\section*{Solution}
Multiplying both sides of the equation by 3, we get

$$
\begin{aligned}
& 3 y=-2 x+9 \\
& 2 x+3 y=9 \quad \text { Standard Form }
\end{aligned}
$$

\section*{Example 1.3.9}
Write the equation $3 x-4 y=10$ in the slope-intercept form.

\section*{Solution}
Solving for $y$, we get

$$
\begin{aligned}
& -4 y=-3 x+10 \\
& y=\frac{3}{4} x-\frac{5}{2} \quad \text { Standard Form }
\end{aligned}
$$

Finally, we learn a very quick and easy way to write an equation of a line in the standard form. But first we must learn to find the slope of a line in the standard form by inspection.

By solving for $y$, it can easily be shown that the slope of the line $A x+B y=C$ is $-A / B$.\\
The reader should verify this.

\section*{Example 1.3.10}
Find the slope of the following lines, by inspection.\\
a $3 x-5 y=10$\\
b. $2 x+7 y=20$\\
c. $4 x-3 y=8$

\section*{Solution}
a $A=3, B=-5$, therefore, $m=-\frac{3}{-5}=\frac{3}{5}$\\
b. $A=2, B=7$, therefore, $m=-\frac{2}{7}$\\
c. $m=-\frac{4}{-3}=\frac{4}{3}$

Now that we know how to find the slope of a line in the standard form by inspection, our job in finding the equation of a line is going to be easy.

\section*{Example 1.3.11}
Find an equation of the line that passes through $(2,3)$ and has slope - $4 / 5$.

\section*{Solution}
Since the slope of the line is $-4 / 5$, we know that the left side of the equation is $4 x+5 y$, and the partial equation is going to be

$$
4 x+5 y=c
$$

Of course, $c$ can easily be found by substituting for $x$ and $y$.

$$
\begin{aligned}
& 4(2)+5(3)=c \\
& 23=c
\end{aligned}
$$

The desired equation is

$$
4 x+5 y=23
$$

If you use this method often enough, you can do these problems very quickly.\\
We summarize the forms for equations of a line below:\\
Slope Intercept form: $\mathbf{y}=\mathbf{m x}+\mathbf{b}$,\\
where $m=$ slope, $b=y$-intercept\\
Point Slope form: $\mathbf{y}-\mathbf{y}_{1}=\mathbf{m}\left(\mathbf{x}-\mathbf{x}_{1}\right)$,\\
where $m=$ slope, $\left(x_{1}, y_{1}\right)$ is a point on the line\\
Standard form: $\mathbf{A x}+\mathbf{B y}=\mathbf{C}$\\
Horizontal Line: $\mathbf{y}=\mathbf{b}$\\
where $b=y$-intercept\\
Vertical Line: $\mathbf{x}=\mathbf{a}$\\
where $a=x$-intercept\\
This page titled 1.3: Desermining the Equation of a Line is shared under a CC BY 4.0 license and was authored, remixed, and/or curated by Rupinder Sekhon and Roberta Bloom via source content that was edited to the style and standards of the LibreTexts platform.

\subsection*{1: Determining the Equation of a Line (Exercises)}
\section*{SECTION 1.3 PROBLEM SET: DETERMINING THE EQUATION OF A LINE}
Write an equation of the line satisfying the following conditions.\\
Write the equation in the form $\mathrm{y}=\mathrm{mx}+\mathrm{b}$.

\begin{enumerate}
  \item It passes through the point $(3,10)$ and has\\
slope $=2$
  \item It passes through point (4,5) and has m $=0$.
  \item It passes through $(3,5)$ and $(2,-1)$.
  \item It has slope 3 , and its $y$-imtercept equals 2 .
  \item It passes through $(5,-2)$ and $m=2 / 5$.
  \item It passes through $(-5,-3)$ and $(10,0)$.\\
7 ) It passes through $(4,-4)$ and $(5,3)$.
  \item It passes through (7, - 2) ; its $y$-intercept is 5 .
  \item It passes through (2, -5) and its $x$-intercept\\
is 4 .
  \item Its a borizontal line through the point ( $2,-1$ ).
\end{enumerate}

\section*{SECTION 1.3 PROBLEM SET: DETERMINING THE EQUATION OF A LINE}
Write an equation of the line satisfying the following conditions. Write the equation in the form $\mathrm{y}=\mathrm{mx}+\mathrm{b}$.\\
11) It passes through ( $5,-4$ ) and ( $1,-4$ ).\\
12) It is a vertical line through the point (3, - 2).\\
13) It passes through $(3,-4)$ and $(3,4)$.\\
14) It has $x$-intercept $=3$ and $y$-intercept $=4$.

Write an equation of the line satisfying the following conditions. Write the equation in the form $\mathrm{Ax}+\mathrm{By}=\mathrm{C}$.\\
15) It passes through (3, -1) and $m=2$.\\
16) It passes through $(-2,1)$ and $m=-3 / 2$.\\
17) It passes through $(-4,-2)$ and $m=3 \times 4$.\\
18) Its $x$-Intercept equals 3 , and $\mathrm{m}=-5 / 3$.

\section*{SECTION 1.3 PROBLEM SET: DETERMINING THE EQUATION OF A LINE}
Write an equation of the line satisfying the following conditions. Write the equation in the form $\mathrm{Ax}+\mathrm{By}=\mathrm{C}$.\\
19) It passes through (2, -3 ) and (5, 1).\\
20) It passes through $(1,-3)$ and $(-5,5)$.

Write an equation of the line satisfying the following conditions.\\
Write the equation in point slope form $y-y_{1}=m\left(x-x_{1}\right)$\\
21) It passes through $(2,-3)$ and $(5,1)$.\\
22) It passes through (1, - 3) and (-5, 2).\\
23) It passes through ( $6,-2$ ) and ( 0,2 ).\\
24) It passes through $(8,2)$ and $(-7,-4)$.\\
25) It passes through $(-12,7)$ and has slope $=-1 / 3$.\\
26) It passes through (8, -7) and has slope 3/4.

\footnotetext{This page titled 1.3.1: Determining the Equation of a Line (Exercises) is shared under a CC BY 4.0 license and was authored, remixed, and/or curated by Rupinder Sekhon and Roberta Bloom via source content that was edited to the style and standark of the LibreTexts platform.
}\section*{4: Applications}
\section*{Learning Objectives}
In this section, you will leam to use linear functions to model real-world applications\\
Now that we have learned to determine equations of lines, we get to apply these ideas in a variety of real-life situations.\\
Read the problem carefully. Highlight important information. Keep track of which values correspond to the independent variable (x) and which correspond to the dependent variable (y).

\section*{Example 1.4.1}
A taxi service charges $\$ 0.50$ per mile plus a $\$ 5$ flat fee. What will be the cost of traveling 20 miles? What will be cost of traveling $x$ miles?

\section*{Solution}
$x=$ distance traveled, in miles and $y=$ cost in dollars\\
The cost of traveling 20 miles is

$$
y=(0.50)(20)+5=10+5=15
$$

The cost of traveling $x$ miles is

$$
y=(0.50)(x)+5=0.50 x+5
$$

In this problem, $\mathbf{S 0 . 5 0}$ per mile is referred to as the variable cost, and the flat charge $\$ 5$ as the fixed cost. Now if we look at our cost equation $y=.50 x+5$, we can see that the variable cost corresponds to the slope and the fixed cost to the $y$-intercept.

\section*{Example 1.4.2}
The variable cost to manufacture a product is $\$ 10$ per item and the fixed cost $\$ 2500$. If $x$ represents the number of items manufactured and $y$ represents the total cost, write the cost function.

\section*{Solution}
\begin{itemize}
  \item The variable cost of $\$ 10$ per item tells us that $m=10$.
  \item The fixed cost represents the $y$-intercept So $b=2500$.
\end{itemize}

Therefore, the cost equation is $y=10 x+2500$.

\section*{Example 1.4.3}
It costs $\$ 750$ to manufacture 25 items, and $\$ 1000$ to manufacture 50 items. Assuming a linear relationship holds, find the cost equation, and use this function to predict the cost of 100 items.

\section*{Solution}
We let $x=$ the number of items manufactured, and let $y=$ the cost.\\
Solving this problem is equivalent to finding an equation of a line that passes through the points ( 25,750 ) and ( 50,1000$)$.

$$
m=\frac{1000-750}{50-25}=10
$$

Therefore, the partial equation is $y=10 x+b$\\
By substituting one of the points in the equation, we get $b=500$\\
Therefore, the cost equation is $y=10 x+500$\\
To find the cost of 100 items, substitute $x=100$ in the equation $y=10 x+500$

So the cost is

$$
y=10(100)+500=1500
$$

It costs $\$ 1500$ to manufacture 100 items.

\section*{Example 1.4.4}
The freezing temperature of water in Celsius is 0 degrees and in Fahrenheit 32 degrees. And the boiling temperatures of water in Celslus, and Fahrenhelt are 100 degrees, and 212 degrees, respectively. Write a conversion equation from Celsius to Fahrenheit and use this equation to convert 30 degrees Celsius into Fahrenhelt.

\section*{Solution}
Let us look at what is given.

\begin{center}
\begin{tabular}{|c|c|}
\hline
Celsius & Fahrenheit \\
\hline
0 & 32 \\
\hline
100 & 212 \\
\hline
\end{tabular}
\end{center}

Again, solving this problem is equivalent to finding an equation of a line that passes through the points ( 0,32 ) and (100, 212).\\
Since we are finding a linear relationship, we are looking for an equation $y=m x+b$, or in this case $F=m C+b$, where $x$ or $C$ represent the temperature in Celsius, and y or F the temperature in Fahrenheit.

$$
\text { slope } m=\frac{312-32}{100-0}=\frac{9}{5}
$$

The equation is $F=\frac{9}{5} C+b$\\
Substituting the point $(0,32)$, we get

$$
F=\frac{9}{5} C+32
$$

To convert 30 degrees Celsius into Fahrenhelt, substitute $C=30$ in the equation

$$
\begin{aligned}
\mathrm{F} & =\frac{9}{5} \mathrm{C}+32 \\
\mathrm{~F} & =\frac{9}{5}(30)+32=86
\end{aligned}
$$

\section*{Example 1.4.5}
The population of Canada in the year 1980 was 24.5 million, and in the year 2010 it was 34 million. The population of Canada over that time period can be approximately modelled by a linear function. Let x represent time as the number of years after 1980 and let y represent the size of the population.\\
a Write the linear function that gives a relationship between the time and the population.\\
b. Assuming the population continues to grow linearly in the future, use this equation to predict the population of Canada in the year 2025.

\section*{Solution}
The problem can be made easier by using 1980 as the base year, that is, we choose the year 1980 as the year zero. This will mean that the year 2010 will correspond to year 30 . Now we look at the information we have:

\begin{center}
\begin{tabular}{|c|c|}
\hline
Year & Population \\
\hline
$0(1980)$ & 24.5 million \\
\hline
$30(2010)$ & 34 million \\
\hline
\end{tabular}
\end{center}

a. Solving this problem is equivalent to finding an equation of a line that passes through the points ( $0,24.5$ ) and (30, 34). We use these two points to find the slope:

$$
m=\frac{34-24.5}{30-0}=\frac{9.5}{30}=0.32
$$

The $y$-intercept occurs when $x=0$, so $b=24.5$

$$
y=0.32 x+24.5
$$

b. Now to predict the population in the year 2025, we let $x=2025-1980=45$

$$
\begin{aligned}
& y=0.32 x+24.5 \\
& y=0.32(45)+24.5=38.9
\end{aligned}
$$

In the year 2025, we predict that the population of Canada will be 38.9 million people.\\
Note that we assumed the population trend will continue to be linear. Therefore if population trends change and this assumption does not continue to be true in the future, this prediction may not be accurate.

\footnotetext{This page titled 1.4: Applications is shared under a CC BY 4.0 license and was authored, remixed, and/or curated by Rupinder Sekhon and Roberta Bloom via source content that was edited to the style and standards of the LibreTexts platform.
}\subsection*{1: Applications (Exercises)}
\section*{SECTION 1.4 PROBLEM SET: APPLICATIONS}
In the following application problems, assume a linear relationship holds.

\textbackslash begin\{abstract\}\\
1) The variable cost to manufacture a product is $\$ 25$ per item, and the fixed costs are $\$ 1200$.\\
If $x$ is the number of items manufactured and $y$ is the cost, write the cost function.\\
3) The variable cost to manufacture an item is\\
$\$ 20$ per item, and it costs a total of $\$ 750$ to produce 20 items. If $x$ represents the number of items manufactured and $y$ is the cost, write the cost function.\\
5) To manutecture 100 items, it costs $\$ 32,000$, and to manufacture 200 items, it costs $\$ 40,000$. If x is the number of items manufactured and $y$ is the cost, write the cost function.\\
\textbackslash end\{abstract\}

\begin{enumerate}
  \setcounter{enumi}{1}
  \item It costs $\$ 90$ to rent a car driven 100 miles and $\$ 140$ for obe driven 200 miles. If x is the number of miles driven and y the total cost of the rental, write the cost function.
  \item To manutacture 30 items, it cosss $\$ 2700$, and to manufacture 50 items, it costs $\$ 2200$. If $x$ represents the number of items manufactured and $y$ the cost, write the cost function.
  \item It costs $\$ 1900$ to manufacture 60 items, and the fixed costs are $\$ 700$. If x represents the number of items manufactured and $y$ the cost, write the cost function.
\end{enumerate}

\section*{SECTION 1.4 PROBLEM SET: APPLICATIONS}
In the following application problems, assume a linear relationship holds.\\
7) A person who weighs 150 pounds has 60 pounds of muscles; a person that weighs 180 pounds has 72 pounds of muscles. If x represents body weight and $y$ is muscle weight, write an equation describing their relationship. Use this relationship to determine the muscle weight of a person that weighs 170 pounds.\\
9). A male college student who is 64 inches tall weighs 110 pounds. Anocher student who is 74 inches tall weighs 180 pounds. Assuming the relationship between male students' heights ( x ), and weights ( y ) is linear, write a function to express weights in terms of heights, and use this function to prediet the weight of a student who is 68 inches tall.\\
11) The freezing temperatures for water for Celsius and Fahrenheit scales are $0^{\circ} \mathrm{C}$ and $32^{\circ} \mathrm{F}$. The boiling temperatures for water are $100^{\circ} \mathrm{C}$ and $212^{\circ}$. Let C denore the temperature in Celsius and F in Fahrenheit. Write the coenversion function from Celsius to Fahrenheit. Use the function to convert $25^{\circ} \mathrm{C}$ into ${ }^{\circ} \mathrm{F}$.\\
B) A spring on a door stretches 6 inches if a force of 30 pounds is applied. It strenches 10 inches If a 50 pound force is applied. If $\times$ represents the number of inches stretched, and $y$ is the farce, write a linear equation describing the relationship. Use it to determine the amount of force required to stretch the spring 12 inches.\\
10) EZ Clean company has determined that if it spends $\$ 30,000$ on advertising, it can hope to sell 12,000 of its Minivacs a year, but if it spends $\$ 50,000$, it can sell 16,000 . Write an equation that gives a relationship between the number of dollars spent on advertising (x) and the number of minivacs sold(y).\\
12) By reversing the coordinates in the previous problem, find a conversion function that convers Fahrenheit inno Celsius, and use this conversion function to convert $72{ }^{\circ} \mathrm{F}$ intn an equivalent Celsius measure.

\section*{SECTION 1.4 PROBLEM SET: APPLICATIONS}
In the following application problems, assume a linear relationship holds.\\
13) California's population was 29.8 million in the year 1990 , and 37.3 million in 2010. Assume that the population trend was and continues to be linear, write the population function. Use this fanction to predict the population in 2025. Hint: Use 1990 as the base year (year 0); then 2010 and 2025 are years 20 , and 35 , respectively.)\\
15) A college's enrollment was 13,200 students in the year 2000, and 15,000 studens in 2015. Enrollment has followed a linear pautur.\\
Write the function that models enrollment as a function of time. Use the function to find the college's enrollment in the year 2010.\\
Hint: Use year 2000 as the base year.\\
14) Use the popalation function for Califomia in the previous problem to find the year in which the population will be 40 million people.\\
16) If the college's enrollment continues to follow this pattern, in what year will the college have 16,000 students enrolled.\\
17) The cost of electricity in residential homes is a linear function of the amount of energy used. In Grove City, a home using 250 kilowat hours (kwh) of electricity per month pays $\$ 55$.\\
A home using 600 kwh per month pays $\$ 118$. Write the cost of electricity as a function of the amount used. Use the function to find the cost for a home using 400 kwh of electricity per monch.\\
18) Find the level of electricity use that would correspond to a monthly cost of \$100.

\section*{SECTION 1.4 PROBLEM SET: APPLICATIONS}
In the following application problems, assume a linear relationship holds.\\
19) At ABC C0., sales revenue is $\$ 170,000$ when it spends $\$ 5000$ on advertising.\\
Sales revenue is $\$ 2 S 4,000$ when $\$ 12,000$ is spent on advertising.\\
a) Find a linear function for\\
$y=$ amount of sales revenue as a function of\\
$\mathrm{x}=$ amount spent on advertising.\\
b) Find revenue if $\$ 10,000$ is spent on advertising.\\
c) Find the amount that should be spent on advertising to achieve $\$ 200,000$ in revenue.\\
21) Mugs Caté sells 1000 cups of coffee per week if it does not advertise. For every $\$ 50$ spent in advertising per week, it sells an additional 150 cups of coffee.\\
a) Find a linear function that gives $y=$ number of cups of cotfee sold per week\\
$x=$ amount spent on advertising per week.\\
b) How many cups of coffee does Mugs Cafe expect to sell if $\$ 100$ per week is spent on advertising?

This page titled 1.4.1: Applications (Exercises) is shared under a CC BY 4.0 license and was authored, remixed, and/or curated by Rupinder Sekhon and Roberta Bloom via source content that was edited to the style and standards of the LibreTexts platform.

\section*{5: More Applications}
\section*{Learning Objectives}
In this section, you will leam to:

\begin{enumerate}
  \item Solve a linear system in two variables.
  \item Find the equilibrium point when a demand and a supply equation are given.
  \item Find the break-even point when the revenue and the cost functions are given.
\end{enumerate}

\section*{Finding the Point of Intersection of Two Lines}
In this section, we will do application problems that involve the intersection of lines. Therefore, before we proceed any further, we will first learn bow to find the intersection of two lines.

\section*{Example 1.5.1}
Find the intersection of the line $y=3 x-1$ and the line $y=-x+7$.

\section*{Solution}
We graph both lines on the same axes, as shown below, and read the solution ( 2,5 ).\\
\includegraphics[max width=\textwidth, center]{2024_12_11_69ed6e5e21c55666356eg-22}

Finding an intersection of two lines graphically is not always easy or practical; therefore, we will now learn to solve these problems algebraically.\\
At the point where two lines intersect, the $x$ and $y$ values for both lines are the same. So in order to find the intersection, we either let the $x$-values or the $y$-values equal.\\
If we were to solve the above example algebraically, it will be easier to let the $y$-values equal. Since $y=3 x-1$ for the first line, and $y=-x+7$ for the second line, by letting the $y$-values equal, we get

$$
\begin{aligned}
3 x-1 & =-x+7 \\
4 x & =8 \\
x & =2
\end{aligned}
$$

By substituting $x=2$ in any of the two equations, we obtain $y=5$.\\
Hence, the solution $(2,5)$.

A common algebraic method used to solve systems of equations is called the elimination method. The object is to eliminate one of the two variables by adding the left and right sides of the equations together. Once one variable is eliminated, we have an equation with only one variable for can be solved. Finally, by substituting the value of the variable that has been found in one of the original equations, we get the value of the other variable.

\section*{Example 1.5.2}
Find the intersection of the lines $2 x+y=7$ and $3 x-y=3$ by the ellimination method.

\section*{Solution}
We add the left and right sides of the two equations.

$$
\begin{aligned}
2 x+y & =7 \\
3 x-y & =3 \\
\hline 5 x & =10 \\
x & =2
\end{aligned}
$$

Now we substitute $x=2$ in any of the two equations and solve for $y$.

$$
\begin{aligned}
2(2)+y & =7 \\
y & =3
\end{aligned}
$$

Therefore, the solution is $(2,3)$.

\section*{Example 1.5.3}
Solve the system of equations $x+2 y=3$ and $2 x+3 y=4$ by the elimination method.

\section*{Solution}
If we add the two equations, none of the variables are eliminated. But the variable $x$ can be eliminated by multiplying the first equation by -2 , and leaving the second equation unchanged.

$$
\begin{aligned}
-2 x-4 y & =-6 \\
2 x+3 y & =4 \\
\hline-y & =-2 \\
y & =2
\end{aligned}
$$

Substituting $y=2$ in $x+2 y=3$, we get


\begin{align*}
& x+2(2)=3  \tag{1.5.1}\\
& x=-1
\end{align*}


Therefore, the solution is $(-1,2)$.

\section*{Example 1.5.4}
Solve the system of equations $3 x-4 y=5$ and $4 x-5 y=6$.

\section*{Solution}
This time, we multiply the first equation by - 4 and the second by 3 before adding. (The choice of numbers is not unique.)

$$
\begin{aligned}
-12 x+16 y & =-20 \\
12 x-15 y & =18 \\
\hline y & =-2
\end{aligned}
$$

By substituting $y=-2$ in any one of the equations, we get $x=-1$.\\
Hence the solution is $(-1,-2)$.

\section*{SUPPLY, DEMAND AND THE EQUILIBRIUM MARKET PRICE}
In a free market economy the supply curve for a commodity is the number of items of a product that can be made available at different prices, and the demand curve is the number of items the consumer will buy at different prices.

As the price of a product increases, its demand decreases and supply increases. On the other hand, as the price decreases the demand increases and supply decreases. The equilibrium price is reached when the demand equals the supply.

\section*{Example 1.5.5}
The supply curve for a product is $y=3.5 x-14$ and the demand curve for the same product is $y=-2.5 x+34$, where x is the price and $y$ the number of items produced. Find the following.\\
a. How many items will be supplied at a price of $\$ 10$ ?\\
b. How many items will be demanded at a price of $\$ 10$ ?\\
c. Determine the equilibrium price.\\
d. How many items will be produced at the equilibrium price?

\section*{Solution}
a) We substitute $x=10$ in the supply equation, $y=3.5 x-14$; the answer is $y=3.5(10)-14=21$.\\
b) We substitute $x=10$ in the demand equation, $y=-2.5 x+34$; the answer is $y=-2.5(10)+34=9$.\\
c) By letting the supply equal the demand, we get

$$
\begin{aligned}
3.5 x-14 & =-25 x+34 \\
6 x & =48 \\
x & =\$ 8
\end{aligned}
$$

d) We substitute $x=8$ in either the supply or the demand equation; we get $y=14$.\\
\includegraphics[max width=\textwidth, center]{2024_12_11_69ed6e5e21c55666356eg-24}

The graph shows the intersection of the supply and the demand functions and their point of intersection, (8,14).\\
Interpretation: At equilibrium, the price is $\$ 8$ per item, and 14 items are produced by suppliers and purchased by consumers.

\section*{The Break-Even Point}
In a business, the profit is generated by selling products.

\begin{itemize}
  \item If a company sells x number of items at a price $P$, then the revenue $\mathbf{R}$ is the price multiplied by number of items sold: $R=P \cdot x$.
  \item The production costs $C$ are the sum of the variable costs and the fixed costs, and are often written as $C=m x+b$, where $x$ is the number of items manufactured.
  \item o The slope $m$ is the called marginal cost and represents the cost to produce one additional item or unit.
  \item The variable cost, mx , depends on how much is being produced
  \item The fixed cost b is constant; it does not change no matter how much is produced.
  \item Profit is equal to Revenue minus Cost: Profit = R - C
\end{itemize}

A company makes a profit if the revenue is greater than the cost. There is a loss if the cost is greater than the revenue. The point on the graph where the revenue equals the cost is called the break-even point. At the break-even point, profit is 0 .

\section*{Example 1.5.6}
If the revenue function of a product is $R=5 x$ and the cost function is $y=3 x+12$, find the following.\\
a If 4 items are produced, what will the revenue be?\\
b. What is the cost of producing 4 items?\\
c. How many items should be produced to break even?\\
d. What will be the revenue and the cost at the break-even point?

\section*{Solution}
a) We substitute $x=4$ in the revenue equation $R=5 x$, and the answer is $R=20$.\\
b) We substitute $x=4$ in the cost equation $C=3 x+12$, and the answer is $C=24$.\\
c) By letting the revenue equal the cost, we get

$$
\begin{aligned}
5 x & =3 x+12 \\
x & =6
\end{aligned}
$$

d) Substitute $x=6$ in either the revenue or the cost equation: we get $R=C=30$.

The graph below shows the intersection of the revenue and cost functions and their point of intersection, (6, 30).\\
\includegraphics[max width=\textwidth, center]{2024_12_11_69ed6e5e21c55666356eg-25}

This page titled 1.5: More Applications is shared under a CC BY 4.0 license and was authored, remixed, and/or curated by Rupinder Sekhon and Roberta Bloom via source content that was edited to the style and standards of the LibreTexts platform.

\subsection*{1: More Applications (Exercises)}
\section*{SECTION 1.5 PROBLEM SET: MORE APPLICATIONS}
\section*{Solve the following problems.}
\begin{enumerate}
  \item Solve for $x$ and $y$.\\
$y=3 x+4$\\
$y=5 x-2$
  \item The supply and demand curves for a product are: Supply $y=2000 x$ -\\
6500\\
Demand $y=-1000 x+28000$,\\
where $x$ is price and $y$ is the number of items. At what price will supply\\
equal demand and how many items will be produced at that price?
\end{enumerate}

$$
\begin{aligned}
& \text { 2) Solve for } x \text { and } y \text {. } \\
& \begin{array}{l}
\text { 2x-3y }=4 \\
3 x-4 y=5
\end{array}
\end{aligned}
$$

\begin{enumerate}
  \setcounter{enumi}{3}
  \item The supply and demand curves for a product are
\end{enumerate}

Supply y = 300x - 18000 and\\
Demand $\mathrm{y}=-100 \mathrm{x}+14000$,\\
where $x$ is price and $y$ is the number of items. At what price will supply equal demand, and how many items will be produced at that price?\\
5) A car rental company offers two plans for one way rentals.

Plan I charges $\$ 36$ per day and 17 cents per mile. Plan II charges $\$ 24$ per day and 25 cents per mile.\\
a. If you were to drive 300 miles in a day, which plan is better?\\
b. For what mileage are both rates equal?

\section*{SECTION 1.5 PROBLEM SET: MORE APPLICATIONS}
Solve the following problems.\\
6) A demand curve for a product is the number of items the consumer will buy at different prices. At a price of $\$ 2$ a store can sell 2400 of a particular type of toy truck. At a price of $\$ 8$ the store can sell 600 such trucks. If $x$ represents the price of trucks and $y$ the number of items sold, write an equation for the demand curve.\\
8) The equilibrium price is the price where the supply equals the demand. From the demand and supply curves obtained in the previous two problems, find the equilibrium price, and determine the number of items that can be sold at that price.

\section*{SECTION 1.5 PROBLEM SET: MORE APPLICATIONS}
Solve the following problems.\\
10) A company's revenue and cost in dollars are given by $\mathrm{R}=225 \mathrm{x}$ and $C=75 x+6000$, where $x$ is the number of items. Find the number of items that must be produced to break-even.\\
12) Whackemhard Sports is planning to introduce a new line of tennis rackets. The fixed costs for the new line are $\$ 25,000$ and the variable cost of producing each racket is $\$ 60$.\\
x is the number of rackets; y is in dollars.\\
If the racket sells for $\$ 80$, how many rackets must be sold in order to break even?\\
7) A supply curve for a product is the oumber of items that can be made available at different prices. A manufacturer of toy trucks can supply 2000 trucks if they are sold for $\$ 8$ each; it can supply only 400 trucks if they are sold for $\$ 4$ each. If $x$ is the price and $y$ the number of items, write an equation for the supply curve.\\
9) A break-even point is the intersection of the cost function and the revenue function, that is, where total cost equals revenue, and profit is zero. Mrs. Jones Cookies Store's cost and revenove, in dollars, for $x$ number of cookies is given by $\mathrm{C}=.05 \mathrm{x}+3000$ and $\mathrm{R}=.80 \mathrm{x}$. Find the number of cookies that must be sold to break even.\\
11) A firm producing socks has a fixed cost of\\
$\$ 20,000$ and variable cost of $\$ 2$ per pair of socks. Let $x=$ the number of pairs of socks. Find the break-even point if the socks sell for $\$ 4.50$ per pair.\\
13) It costs $\$ 1,200$ to produce 50 pounds of a chemical and it costs $\$ 2,200$ to produce 150 pounds. The chemical sells for $\$ 15$ per pound x is the amount of chemical; y is in dollars.\\
a. Find the cost function.\\
b. What is the fixed cost?\\
c How many pounds must be sold to break even?\\
d. Find the cost and revenue at the break-even point.

\footnotetext{This page titled 1.5.1: More Applications (Exercises) is shared under a CC BY 4.0 license and was authored, remixed, and/or curated by Rupinder
}Sekhon and Roberta Bloom via source content that was edited to the style and standards of the LibreTexts platform.

\section*{6: Chapter Review}
\section*{SECTION 1.6 PROBLEM SET: CHAPTER REVIEW}
\begin{enumerate}
  \item Find an equation of the $x$-axis.
  \item Find the slope of the line whose equation is $2 x+3 y=6$.
\end{enumerate}

3 ) Find the slope of the line whose equation is $y=-3 x+5$.\\
4) Find both the $x$ and $y$ intercepts of the line $3 x-2 y=12$.\\
5) Find an equation of the line whose slope is 3 and $y$-intercept 5.\\
6) Find an equation of the line whose $x$-intercept is 2 and $y$-intercept 3 .\\
7) Find an equation of the line that has slope 3 and passes through the point (2, 15).\\
8) Find an equation of the line that has slope $-3 / 2$ and passes through the point $(4,3)$.\\
9) Find an equation of the line that passes through the points $(0,32)$ and $(100,212)$.\\
10) Find an equation of the line that passes through the point $(2,5)$ and is parallel to the line $y=3 x+4$.\\
11) Find the point of intersection of the lines $2 x-3 y=9$ and $3 x+4 y=5$.\\
12) Is the point $(3,-2)$ on the line $5 x-2 y=11$ ?\\
13) Find two points on the line given by the parametric equations, $x=2+3 \mathrm{t}, \mathrm{y}=1-2 \mathrm{t}$.\\
14) Find two points on the line $2 x-6=0$.\\
15) Graph the line $2 x-3 y+6=0$.\\
16) Graph the line $\mathrm{y}=-2 \mathrm{x}+3$.\\
17) A female college student who is 60 inches tall weighs 100 pounds. Another female student who is 66 inches tall weighs 124 pounds. Assume the relationship between the female students' weights and heights is linear. Find an equation for weight as a function of height. Use this relationship to predict the weight of a female student who is 70 inches tall.\\
18) In deep-sea diving, the pressure exerted by water plays a great role in designing underwater equipment. If at a depth of 10 feet there is a pressure of $21 \mathrm{lb} / \mathrm{In}^{2}$, and at a depth of 50 ft there is a pressure of $75 \mathrm{lb} / \mathrm{in}^{2}$, write a linear equation giving a relationship between depth and pressure. Use this relationship to predict pressure at a depth of 100 ft .\\
19) The variable cost to manufacture an item is $\$ 30$ per item; the fixed costs are $\$ 2750$. Find the cost function.\\
20) The variable cost to manufacture an item is $\$ 10$ per item, and it costs $\$ 2,500$ to produce 100 items. Write the cost function, and use this function to estimate the cost of manufacturing 300 items.\\
21) It costs $\$ 2,700$ to manufacture 100 items of a product, and $\$ 4,200$ to manufacture 200 items. $x=$ the number of items; $y=$ cost. Find the cost function; use it to predict the cost to produce 1000 items.\\
22) In 1990, the average bouse in Emerald City cost $\$ 280,000$ and in 2007 the same house cost $\$ 365,000$. Assuming a linear relationship, write an equation that will give the price of the house in any year, and use this equation to predict the price of a similar house in the year 2020.\\
23) The population of Mexico in 1995 was 95.4 million and in 2010 it was 117.9 million. Assuming a linear relationship, write an equation that will give the population of Mexico in any year, and use this equation to predict the population of Mexico in the year 2025.

\section*{SECTION 1.6 PROBLEM SET: CHAPTER REVIEW}
\begin{enumerate}
  \setcounter{enumi}{23}
  \item At Nuts for Soup Lunch Bar, they sell 150 bowls of soup if the high temperature for the day is $40^{\circ} \mathrm{F}$. For every $5^{\circ} \mathrm{F}$ increase in high temperature for the day, they sell 10 fewer bowls of soup.\\
a Assuming a linear relationship, write an equation that will give $y=$ the number of bowls of soup sold as a function of $x=$ the daily high temperature.\\
b. How many bowls of soup are sold when the temperature is $75^{\circ}$ F?\\
c. What is the temperature when 100 bowls of soup are sold?
  \item Two-hundred items are demanded at a price of $\$ 5$, and 300 items are demanded at a price of $\$ 3$. If $x$ represents the price, and $y$ the number of items, write the demand function.
  \item A supply curve for a product is the number of items of the product that can be made available at different prices. A doll manufacturer can supply 2000 dolls if the dolls are sold for $\$ 30$ each, but he can supply only 400 dolls if the dolls are sold for $\$ 10$ each. If $x$ represents the price of dolls and $y$ the number of items, write an equation for the supply curve.
  \item Suppose you are trying to decide on a price for your latest creation - a coffee mug that never tips. Through a survey, you have determined that at a price of $\$ 2$, you can sell 2100 mugs, but at a price of $\$ 12$ you can only sell 100 mugs. Furthermore, your supplier can supply you 3100 mugs if you charge your customers $\$ 12$, but only 100 mugs if you charge $\$ 2$. What price should you charge so that the supply equals demand, and at that price how many coffee mugs will you be able to sell?
  \item A car rental company offers two plans. Plan I charges $\$ 16$ a day and 25 cents a mile, while Plan II charges $\$ 45$ a day but no charge for miles. If you were to drive 200 miles in a day, which plan is better? For what mileage are both rates the same?
  \item The supply curve for a product is $y=250 x-1000$. The demand curve for the same product is $y=-350 x+8,000$, where $x$ is the price and $y$ the number of items produced. Find the following.\\
a At a price of $\$ 10$, how many items will be in demand?\\
b. At what price will 4,000 items be supplied?\\
c. What is the equilibrium price for this product?\\
d. How many items will be manufactured at the equilibrium price?
  \item The supply curve for a product is $y=625 x-600$ and the demand curve for the same product is $y=-125 x+8,400$, where $x$ is the price and $y$ the number of items produced.\\
Find the equilibrium price and determine the number of items that will be produced at that price.
  \item Both Jenny and Masur work in the sales department for Sports Supply. Jenny gets paid $\$ 120$ per day plus $4 \%$ commission on the sales. Masur gets paid $\$ 132$ per day plus $8 \%$ commission on the sales in excess of $\$ 1,000$. For what sales amount would they both eam the same daily amounts?
  \item A company's revenue and cost in dollars are given by $\mathrm{R}=25 \mathrm{x}$ and $\mathrm{C}=10 \mathrm{x}+9,000$, where x represents the number of items. Find the number of items that must be produced to break-even.
  \item A firm producing a certain type of CFL lightbulb has fixed costs of $\$ 6,800$, and a variable cost of $\$ 2.30$ per bulb. The bulbs sell for $\$ 4$ each. How many bulbs must be produced to break-even?
  \item A company producing tire pressure gauges has fixed costs of $\$ 7,500$, and variable cost of $\$ 1.50$ cents per item. If the gauges sell for $\$ 4.50$, how many must be produced to break-even?
  \item A company is introducing a new cordless travel shaver before the Christmas holidays. It hopes to sell 15,000 of these shavers in December. The variable cost is $\$ 11$ per item and the fixed costs $\$ 100,000$. If the shavers sell for $\$ 19$ each, how many must be produced and sold to break-even?
\end{enumerate}

\footnotetext{This page titled 1.6: Chapter Review is shared under a CC BY 4.0 livense and was authored, remixed, and/or curated by Rupinder Sekblon and Roberta Bloom via source content that was edited to the style and standards of the LibreTexts platform.
}
